\documentclass[10pt,a4paper]{amsart}
\usepackage[margin=19mm]{geometry}
\usepackage{amsmath,amssymb,mathtools}
\usepackage[scr=rsfs,cal=euler]{mathalfa}
\usepackage{xfrac}
\usepackage[unicode,hidelinks,bookmarks=false]{hyperref}
\newcommand{\ie}{i.e.\ }
\newcommand{\cf}{cf.\ }
\newcommand{\resp}{respectively\ }
\newcommand{\Subsection}{\subsection}
\providecommand{\nobreakdash}{\nobreak\hbox{-}}
\setlength{\emergencystretch}{2em}
\multlinegap0pt
\usepackage{amssymb}
\newtheorem{prop}{Proposition}
\DeclareMathOperator{\SYM}{SYM}
\DeclareMathOperator{\Y}{\mathbb{Y}}
\title{Probability measures on families of partitions related to harmonic analysis on big wreath products}
\author{Eugene Strahov}
\date{Source: arXiv:2505.08394v1 (CC BY 4.0)}
\begin{document}
\maketitle

\noindent\emph{Source and licence.} Probability measures on families of partitions related to harmonic analysis on big wreath products. Version arXiv:2505.08394v1 (CC BY 4.0), distributed by arXiv under the Creative Commons Attribution 4.0 International licence, http://creativecommons.org/licenses/by/4.0/, as recorded in the arXiv licence metadata for this exact version and shown on the abstract page https://arxiv.org/abs/2505.08394v1. Full licence notice: \texttt{licenses/arxiv-2505.08394-CC-BY-4.0.txt}.\par\smallskip

\noindent\textit{Editorial scope.} Counted unit: the article's prop PropPT6 (source0.tex lines 1642--1650). The article's complete original proof of it is delivered with it (source0.tex lines 1651--1727, one \texttt{proof} environment of 77 lines). No mathematics is written, repaired or re-derived: the statement and the proof are the article's own lines, in the article's own notation, and no retained line is substituted or re-flowed.  Retained uncounted as the source-local context the proof needs: lines 409--463; lines 1053--1055; lines 1367--1370; lines 1414--1435; lines 1422--1429; lines 1635--1640.  Nothing was omitted from any retained span, so the package carries no omission record.  No retained line contains a citation, so the wrapper carries no bibliography and no bibliography entry.  No retained line contains a figure environment, an image-inclusion command or drawing code, so no image is delivered and the package declares no asset dependency.  Editorial formatting only, in addition: an AMS \texttt{amsart} preamble that restates the article's own declarations and macros used by the retained lines, namely the environments \texttt{prop} on separate counters, together with the standard packages those lines need.  The retained environments print in this document as Proposition 1 in order of appearance, whereas the article prints the same items with the numbering of its own full document.  The complete native source of the article as distributed in the arXiv e-print for this exact version is bundled unchanged as \texttt{sources/178038/source0.tex}.
\begin{equation}
D_n(G)=G^n=\left\{d=(g_1,\ldots,g_n)\vert g_1\in G,\ldots,g_n\in G\right\}.
\end{equation}
The action of $S_n$ on $D_n(G)=G^n$ is defined by
$$
s:\; d=\left(g_1,\ldots,g_n\right)\longrightarrow s(d)=\left(g_{s^{-1}(1)},\ldots,g_{s^{-1}(n)}\right).
$$
Denote by $S_n(G)$ the wreath product of $G$ with $S_n$. By definition, $S_n(G)$ is the semidirect product $S_n(G)=D_n(G)\rtimes S_n$ (that is, each element of $S_n(G)$ has the form
$(d,s)$, $d\in D_n(G)$, $s\in S_n$, and $(d,s)(d',s')=(ds(d'),ss')$).
In other words,  the wreath product $S_n\left(G\right)$
is a group whose underlying set is
$$
G^n\times S_n=\left\{\left(\left(g_1,\ldots,g_n\right),s\right):\; g_i\in G, s\in S_n\right\}.
$$
The multiplication in $S_n\left(G\right)$ can be expressed explicitly as
$$
\left(\left(g_1,\ldots,g_n\right),s\right)
\left(\left(h_1,\ldots,h_n\right),t\right)
=\left(\left(g_1h_{s^{-1}(1)},\ldots,g_nh_{s^{-1}(n)}\right),st\right).
$$
When $n=1$, $S_1\left(G\right)$ is $G$.
%%%%%%%%%%%%%%%%%%%%%%%%%%%%%%%%%%%%%%%%%%%%%%%%%%%%%%%%%%%%%%%%%%%%%%%%%%%%%%%%%%
%%%%%%%%%%%%%%%%%%%%%%%%%%%%%%%%%%%%%%%%%%%%%%%%%%%%%%%%%%%%%%%%%%%%%%%%%%%%%%%%%%%%%%
%%%%%%%%%%%%%%%%%%%%%%%%%%%%%%%%%%%%%%%%%%%%%%%%%%%%%%%%%%%%%%%%%%%%%%%%%%%%%%%%%%%%%
%%%%%%%%%%%%%%%%%%%%%%%%%%%%%%%%%%%%%%%%%%%%%%%%%%%%%%%%%%%%%%%%%
%%%%%%%%%%%%%%%%%%%%%%%%%%%%%%%%%%%%%%%%%%%%%%%%%%%%%%%%%%%%%%%%%%%%%%%%%%
\subsection{Conjugacy classes and irreducible representations of $G$}\label{NotSection1.2}
%%%%%%%%%%%%%%%%%%%%%%%%%%%%%%%%%%%%%%%%%%%%%%%%%%%%%%%%%%%%%%%%%%%%%%%%%%%%%%%%%%%%%%%%%%%%%%%%%%
%%%%%%%%%%%%%%%%%%%%%%%%%%%%%%%%%%%%%%%%%%%%%%%%%%%%%%%%%%%%%%%%%%%%%%%%%%%%%%%%%%
%%%%%%%%%%%%%%%%%%%%%%%%%%%%%%%%%%%%%%%%%%%%%%%%%%%%%%%%%%%%%%%%%%%%%%%%%%%%%%%%%%%%%%
%%%%%%%%%%%%%%%%%%%%%%%%%%%%%%%%%%%%%%%%%%%%%%%%%%%%%%%%%%%%%%%%%%%%%%%%%%%%%%%%%%%%%
%%%%%%%%%%%%%%%%%%%%%%%%%%%%%%%%%%%%%%%%%%%%%%%%%%%%%%%%%%%%%%%%%
Denote by $\left[G\right]$ the set of conjugacy classes
of $G$, and by $\widehat{G}$ the set of equivalence classes of continuous irreducible
unitary representations of $G$. We assume that $\widehat{G}$ is at most countable.
Therefore, an element $\zeta\in\widehat{G}$ is an equivalent class of continuous unitary irreducible representations of $G$.  The notation $\pi^{\zeta}\in\zeta$ means that $\pi^{\zeta}$ is a representation of class $\zeta$. We denote by $V^{\zeta}$ the representation space
of $\pi^{\zeta}$, and by $\chi^{\pi^\zeta}$ the character of $\pi^{\zeta}$.
%%%%%%%%%%%%%%%%%%%%%%%%%%%%%%%%%%%%%%%%%%%%%%%%%%%%%%%%%%%%%%%%%%%%%%%%%%%%%%%%%%
%%%%%%%%%%%%%%%%%%%%%%%%%%%%%%%%%%%%%%%%%%%%%%%%%%%%%%%%%%%%%%%%%%%%%%%%%%%%%%%%%%%%%%
%%%%%%%%%%%%%%%%%%%%%%%%%%%%%%%%%%%%%%%%%%%%%%%%%%%%%%%%%%%%%%%%%%%%%%%%%%%%%%%%%%%%%
%%%%%%%%%%%%%%%%%%%%%%%%%%%%%%%%%%%%%%%%%%%%%%%%%%%%%%%%%%%%%%%%%
\subsection{Representations of $D_n(G)=G^n$}\label{SubsectionDn}
%%%%%%%%%%%%%%%%%%%%%%%%%%%%%%%%%%%%%%%%%%%%%%%%%%%%%%%%%%%%%%%%%%%%%%%%%%%%%%%%%%
%%%%%%%%%%%%%%%%%%%%%%%%%%%%%%%%%%%%%%%%%%%%%%%%%%%%%%%%%%%%%%%%%%%%%%%%%%%%%%%%%%%%%%
%%%%%%%%%%%%%%%%%%%%%%%%%%%%%%%%%%%%%%%%%%%%%%%%%%%%%%%%%%%%%%%%%%%%%%%%%%%%%%%%%%%%%
%%%%%%%%%%%%%%%%%%%%%%%%%%%%%%%%%%%%%%%%%%%%%%%%%%%%%%%%%%%%%%%%%
Denote by $\widehat{D_n(G)}$ the set of equivalence classes of continuous unitary irreducible representations of $D_n(G)$.  We can represent $\widehat{D_n(G)}$ as
$$
\widehat{D_n(G)}=\left\{\eta=\left(\zeta^1,\ldots,\zeta^n\right)\vert
\zeta^1\in\widehat{G},\ldots,\zeta^n\in\widehat{G}\right\}.
$$
Each representation of $D_n(G)$ from the equivalence class $\eta=\left(\zeta^1,\ldots,\zeta^n\right)$ can be written as
\begin{equation}\label{reps1}
\pi^{\zeta^1}\boxtimes\ldots\boxtimes\pi^{\zeta^n},
\end{equation}

\begin{equation}\label{ch1.3.3}
\widehat{p}_r(\zeta)=\left<p_r,\chi^{\pi^{\zeta}}\right>_G,
\end{equation}

\begin{equation}\label{HHH3.5.1}
\varrho^{\zeta}(g,s)=\left(\pi^{\zeta}\right)^{\boxtimes\; |\lambda^{\zeta}|}
(g)I(s),\;\;(g,s)\in S_{|\lambda^{\zeta}|}(G).
\end{equation}

\begin{prop}\label{chchiProposition5.1.5}
Let $\left(\pi^{\zeta},V^{\zeta}\right)$ denote the irreducible unitary representation of
a compact group $G$ parameterized by $\zeta\in\widehat{G}$. Assume that
\begin{equation}\label{chchi5.1.2}
s=\left(i_1i_2\ldots i_{l_1}\right)\left(j_1j_2\ldots j_{l_2}\right)\ldots\left(\nu_1\nu_2\ldots\nu_{l_p}\right)\in S_{|\lambda^{\zeta}|},
\end{equation}
so $l_1+l_2+\ldots+l_p=|\lambda^{\zeta}|$.
Then we have
\begin{equation}\label{chchi5.1.3}
\begin{split}
&\chi^{\varrho^{\zeta}}\left[\left((g_1,\ldots,g_{|\lambda^{\zeta}|}),s\right)\right]\\
&=\chi^{\pi^{\zeta}}\left(g_{i_{l_1}}\ldots g_{i_2}g_{i_1}\right)
\chi^{\pi^{\zeta}}\left(g_{j_{l_2}}\ldots g_{j_2}g_{j_1}\right)
\ldots \chi^{\pi^{\zeta}}\left(g_{\nu_{l_p}}\ldots g_{\nu_2}g_{\nu_1}\right),
\end{split}
\end{equation}
where $\chi^{\pi^{\zeta}}$ denotes the character of the irreducible representation $\left(\pi^{\zeta},V^{\zeta}\right)$
of $G$ (see Section \ref{NotSection1.2}),
and $\chi^{\varrho^{\zeta}}$ denotes the character of the representation
$\left(\varrho^{\zeta},\left(V^{\zeta}\right)^{\otimes |\lambda^{\zeta}|}\right)$ of the group
$S_{|\lambda^{\zeta}|}(G)$ defined by (\ref{HHH3.5.1}).
\end{prop}

\begin{equation}
\begin{split}
&\chi^{\varrho^{\zeta}}\left[\left((g_1,\ldots,g_{|\lambda^{\zeta}|}),s\right)\right]\\
&=\chi^{\pi^{\zeta}}\left(g_{i_{l_1}}\ldots g_{i_2}g_{i_1}\right)
\chi^{\pi^{\zeta}}\left(g_{j_{l_2}}\ldots g_{j_2}g_{j_1}\right)
\ldots \chi^{\pi^{\zeta}}\left(g_{\nu_{l_p}}\ldots g_{\nu_2}g_{\nu_1}\right),
\end{split}
\end{equation}

\begin{equation}\label{PT5}
J_{\zeta}(s)=\int\limits_{G}\ldots\int\limits_{G}
\Psi\left[\left((g_1,\ldots,g_n),s\right)\right]
\overline{\chi^{\varrho^{\zeta}}\left[\left((g_1,\ldots,g_n),s\right)\right]}
d\mu_G(g_1)\ldots d\mu_{G}(g_n),
\end{equation}

\begin{prop}\label{PropPT6}
Assume that for each $\zeta\in\widehat{G}$, and each $r=1,2,\ldots $ the functions $\widehat{p}_r(\zeta)$ are defined by equation (\ref{ch1.3.3}). Then $J_{\zeta}(s)$ defined by equation (\ref{PT5}) is given by
\begin{equation}\label{PT7}
J_{\zeta}(s)=\widehat{p}_{\mu(s)}(\zeta),\;\; s\in S_n,
\end{equation}
where the Young diagram $\mu(s)$ determines the cyclic structure of $s$, and
$\widehat{p}_{\mu(s)}(\zeta)$ is the Newton power symmetric function in $\widehat{\SYM}(\zeta)$
parameterized by $\mu(s)$.
\end{prop}

\begin{proof}
Recall that the characters $\chi^{\varrho^{\zeta}}$ are given by Proposition \ref{chchiProposition5.1.5}. First, assume that $s=(i_1i_2\ldots i_n)$. From equation
(\ref{chchi5.1.3}) we obtain
\begin{equation}\label{PT7.1}
\chi^{\varrho^{\zeta}}\left[\left((g_1,\ldots,g_n),(i_1i_2\ldots i_n)\right)\right]
=\chi^{\pi^{\zeta}}\left(g_{i_n}\ldots g_{i_2}g_{i_1}\right).
\end{equation}
The conjugacy class of $\left((g_1,\ldots,g_n),(i_1i_2\ldots i_n)\right)$ in $S_n(G)$
is determined by the map $\varrho: [G]\longrightarrow\Y$ defined by
$$
\varrho(c)=\left\{
  \begin{array}{ll}
    (n), & \hbox{$c$ is the conjugacy class of $g_{i_n}\ldots g_{i_2}g_{i_1}$}, \\
    \emptyset, & \hbox{otherwise}
  \end{array}
\right.
$$
This gives
\begin{equation}\label{PT.8}
\Psi\left[\left((g_1,\ldots,g_n),(i_1i_2\ldots i_n)\right)\right]=p_n\left(y_1(g_{i_n}\ldots g_{i_2}g_{i_1}),
y_2(g_{i_n}\ldots g_{i_2}g_{i_1}),\ldots\right).
\end{equation}
Inserting (\ref{PT7.1}), (\ref{PT.8}) into (\ref{PT5}), and
using the invariance of the Haar measure under the shifts by elements of the group, we find
\begin{equation}
\begin{split}
J_{\zeta}\left((i_1i_2\ldots i_n)\right)=\left<p_n\left(y_1(.),y_2(.),\ldots\right),
\chi^{\pi^{\zeta}}(.)\right>_{G}
=\widehat{p}_n(\zeta),
\end{split}
\end{equation}
so equation (\Ref{PT7}) holds true for $s=(i_1i_2\ldots i_n)$.

Next, assume that $s\in S_n$ is given by
$$
s=(i_1i_2\ldots i_l)(j_1j_2\ldots j_k),\;\; k+l=n.
$$
Then (\ref{chchi5.1.3}) enables us to write
\begin{equation}\label{PT9}
\chi^{\varrho^{\zeta}}\left[\left((g_1,\ldots,g_n),(i_1i_2\ldots i_l)(j_1j_2\ldots j_k)\right)\right]=
\chi^{\pi^{\zeta}}(g_{i_l}\ldots g_{i_2}g_{i_1})\chi^{\pi^{\zeta}}(g_{j_k}\ldots g_{j_2}g_{j_1}).
\end{equation}
Denote by $c_1$ the conjugacy class of $g_{i_l}\ldots g_{i_1}g_{i_1}$ in $G$, and
by $c_2$ the conjugacy class of $g_{j_k}\ldots g_{j_2}g_{j_1}$ in $G$. Then the function
$\varrho: [G]\rightarrow \Y$ describing the conjugacy class of
$\left((g_1,\ldots,g_n), (i_1i_2\ldots i_l)(j_1j_2\ldots j_k)\right)$ in $S_n(G)$
can be represented as
$$
\varrho(c)=\left\{
  \begin{array}{ll}
    (l), & c=c_1, c_1\neq c_2, \\
    (k), & c=c_2,c_1\neq c_2, \\
    (l)\cup (k), & c=c_1, c_1=c_2, \\
    \emptyset, & \hbox{otherwise.}
  \end{array}
\right.
$$
This implies the following formula
\begin{equation}\label{PT10}
\begin{split}
&\Psi\left[\left((g_1,\ldots,g_n),(i_1i_2\ldots i_l)(j_1j_2\ldots j_k)\right)\right]\\
&=p_{(l)}\left(y_1(c_1),y_2(c_1),\ldots\right)p_{(k)}\left(y_1(c_2),y_2(c_2),\ldots\right).
\end{split}
\end{equation}
Inserting (\ref{PT9}) and (\ref{PT10}) into (\ref{PT5}) we find
\begin{equation}
\begin{split}
&J_{\zeta}\left((i_1i_2\ldots i_l)(j_1j_2\ldots j_k)\right)\\
&=\left<p_l\left(y_1(.),y_2(.),\ldots\right),
\chi^{\pi^{\zeta}}(.)\right>_{G}\left<p_k\left(y_1(.),y_2(.),\ldots\right),
\chi^{\pi^{\zeta}}(.)\right>_{G}\\
&=\widehat{p}_l(\zeta)\widehat{p}_k(\zeta),
\end{split}
\end{equation}
so (\ref{PT7}) holds true for $s=(i_1i_2\ldots i_l)(j_1j_2\ldots j_k)$ as well. The general case where
$s\in S_n$ is a product of an arbitrary number of cycles is considered in the same way.
\end{proof}

\end{document}
